\section{杨--米尔斯理论与能量--动量张量}
\label{sec:2}
\subsection{维数正规化下的能量--动量张量}
本文考虑定义在 $D$~维欧几里得空间中的 $SU(N)$ 杨--米尔斯理论。作用量为
\begin{equation}
   S=\frac{1}{4g_0^2}\int d^Dx\,F_{\mu\nu}^a(x)F_{\mu\nu}^a(x),
\label{eq:(2.1)}
\end{equation}
其中杨--米尔斯场强为
\begin{equation}
   F_{\mu\nu}(x)
   =\partial_\mu A_\nu(x)-\partial_\nu A_\mu(x)+[A_\mu(x),A_\nu(x)].
\label{eq:(2.2)}
\end{equation}
我们取
\begin{equation}
   D=4-2\epsilon,
\label{eq:(2.3)}
\end{equation}
于是裸规范耦合~$g_0$ 的质量量纲为~$\epsilon$。

假设该理论用维数正规化来正规化（极好的阐述见文
献~\cite{Collins:1984xc}），则系统~\eqref{eq:(2.1)} 的能量--动量张量可以简
单地定义为（例如见文献~\cite{Freedman:1974gs}）
\begin{equation}
   T_{\mu\nu}(x)=\frac{1}{g_0^2}
   \left[F_{\mu\rho}^a(x)F_{\nu\rho}^a(x)
   -\frac{1}{4}\delta_{\mu\nu}F_{\rho\sigma}^a(x)F_{\rho\sigma}^a(x)\right],
\label{eq:(2.4)}
\end{equation}
其中来自规范固定与 Faddeev--Popov 鬼场的项被略去——它们在规范不变算符的关
联函数中不重要。注意，能量--动量张量的质量量纲为~$D$。

维数正规化的优点是其平移不变性。由于这一性质，由裸量朴素地构造的能量--动
量张量、即式~\eqref{eq:(2.4)}，是守恒的，并通过 WT 关系生成正确归一化的平
移：
\begin{equation}
   \int d^Dx\,\left\langle\partial_\mu T_{\mu\nu}(x)\mathcal{O}\right\rangle
   =-\left\langle\partial_\nu\mathcal{O}\right\rangle,
\label{eq:(2.5)}
\end{equation}
这里理解为：右边的导数作用在规范不变算符~$\mathcal{O}$ 中的所有位置上。
把这一 WT 关系与量纲计数及规范不变性结合起来，可知能量--动量张
量~$T_{\mu\nu}(x)$ 是有限的~\cite{Joglekar:1975jm,Nielsen:1977sy}；因而在
最小减除（MS）方案中\footnote{这里我们把重正化算符定义为减去其真空期望值
之后的算符。在使用维数正规化的微扰论中，这一减除是自动的。}有
\begin{equation}
   T_{\mu\nu}(x)-\left\langle T_{\mu\nu}(x)\right\rangle
   =\left\{T_{\mu\nu}\right\}_R(x).
\label{eq:(2.6)}
\end{equation}

能量--动量张量~\eqref{eq:(2.4)} 的有限性为维数~$4$ 规范不变算符的重正化提
供了进一步的有用信息。维数正规化下规范耦合的重正化定义为
\begin{equation}
   g_0^2\equiv\mu^{2\epsilon}g^2Z,
\label{eq:(2.7)}
\end{equation}
其中 $\mu$ 是重正化标度，$Z$ 是重正化因子。在 MS 方案中，
\begin{equation}
   Z=1-\frac{1}{\epsilon}
   \left[b_0g^2+\frac{1}{2}b_1g^4+O(g^6)\right]
   +O\left(\frac{1}{\epsilon^2}\right),
\label{eq:(2.8)}
\end{equation}
以及
\begin{equation}
   b_0=\frac{11N}{48\pi^2},\qquad
   b_1=\frac{17N^2}{384\pi^4}.
\label{eq:(2.9)}
\end{equation}
由维数正规化所保持的转动不变性可以看出，算符重正化具有如下结
构：\footnote{同样，我们定义重正化算符为减去真空期望值之后的算符。}
\begin{align}
   &F_{\mu\rho}^a(x)F_{\nu\rho}^a(x)
   -\left\langle F_{\mu\rho}^a(x)F_{\nu\rho}^a(x)\right\rangle
\notag\\
   &=Z_T\left\{F_{\mu\rho}^aF_{\nu\rho}^a\right\}_R(x)
   +Z_M\delta_{\mu\nu}\left\{F_{\rho\sigma}^aF_{\rho\sigma}^a\right\}_R(x),
\label{eq:(2.10)}
\end{align}
以及
\begin{equation}
   F_{\rho\sigma}^a(x)F_{\rho\sigma}^a(x)
   -\left\langle F_{\rho\sigma}^a(x)F_{\rho\sigma}^a(x)\right\rangle
   =Z_S\left\{F_{\rho\sigma}^aF_{\rho\sigma}^a\right\}_R(x).
\label{eq:(2.11)}
\end{equation}
把上述关系代入式~\eqref{eq:(2.4)} 与~\eqref{eq:(2.6)}，得到
\begin{align}
   &\left\{T_{\mu\nu}\right\}_R(x)
\notag\\
   &=\frac{1}{g^2}\mu^{-2\epsilon}
   Z^{-1}\biggl[Z_T\left\{F_{\mu\rho}^aF_{\nu\rho}^a\right\}_R(x)
   -\frac{1}{4}(Z_S-4Z_M)
   \delta_{\mu\nu}\left\{F_{\rho\sigma}^aF_{\rho\sigma}^a\right\}_R(x)
   \biggr].
\label{eq:(2.12)}
\end{align}
由于左边在~$\epsilon\to0$ 时有限，而在只减除极点项的 MS 方案中，我们推断
（考虑 $\mu\neq\nu$ 与~$\mu=\nu$ 两种情形）：
\begin{equation}
   Z_T=Z=1-b_0g^2\frac{1}{\epsilon}+O(g^4)
\label{eq:(2.13)}
\end{equation}
以及
\begin{equation}
   Z_S-4Z_M=Z.
\label{eq:(2.14)}
\end{equation}

\subsection{迹反常的含义}
能量--动量张量~\eqref{eq:(2.4)} 的另一个重要性质是迹反常~\cite{Adler:1976zt,%
Nielsen:1977sy,Collins:1976yq}：
\begin{equation}
   \delta_{\mu\nu}\left\{T_{\mu\nu}\right\}_R(x)
   =-\frac{\beta}{2g^3}
   \left\{F_{\rho\sigma}^aF_{\rho\sigma}^a\right\}_R(x).
\label{eq:(2.15)}
\end{equation}
由式~\eqref{eq:(2.6)}，此关系等价于
\begin{align}
   \delta_{\mu\nu}
   \left[T_{\mu\nu}(x)-\left\langle T_{\mu\nu}(x)\right\rangle\right]
   &=\epsilon\frac{1}{2g_0^2}
   F_{\rho\sigma}^a(x)F_{\rho\sigma}^a(x)
   -\left\langle\epsilon\frac{1}{2g_0^2}
   F_{\rho\sigma}^a(x)F_{\rho\sigma}^a(x)\right\rangle
\notag\\
   &\xrightarrow{\epsilon\to0}
   -\frac{\beta}{2g^3}
   \left\{F_{\rho\sigma}^aF_{\rho\sigma}^a\right\}_R(x).
\label{eq:(2.16)}
\end{align}
在式~\eqref{eq:(2.15)} 与~\eqref{eq:(2.16)} 中，$\beta$ 表示 $D=4$ 时的
$\beta$ 函数，其定义为
\begin{equation}
   \beta\equiv\left(\mu\frac{\partial}{\partial\mu}\right)_0g
   =-\frac{1}{2}g\left(\mu\frac{\partial}{\partial\mu}\right)_0\ln Z,
\label{eq:(2.17)}
\end{equation}
其中下标~$0$ 表示求导时保持裸量固定。式~\eqref{eq:(2.8)} 与~\eqref{eq:(2.7)}
给出
\begin{equation}
   \beta=-b_0g^3-b_1g^5+O(g^7).
\label{eq:(2.18)}
\end{equation}
然后，把式~\eqref{eq:(2.12)} 代入式~\eqref{eq:(2.15)} 并利用
式~\eqref{eq:(2.13)} 与~\eqref{eq:(2.14)}，我们看到
\begin{equation}
   \delta_{\rho\lambda}
   \left\{F_{\rho\sigma}^aF_{\lambda\sigma}^a\right\}_R(x)
   =\left(1-\frac{\beta}{2g}\right)
   \left\{F_{\rho\sigma}^aF_{\rho\sigma}^a\right\}_R(x),
\label{eq:(2.19)}
\end{equation}
也就是说，与度规的缩并和最小减除（即 $1/\epsilon$ 极点的减除）互不对易；
这是维数正规化的一个奇特但正当的性质~\cite{Collins:1984xc}。

另外，把式~\eqref{eq:(2.7)} 与~\eqref{eq:(2.11)} 代入式~\eqref{eq:(2.16)}，
可得
\begin{equation}
   \epsilon\frac{Z_S}{Z}
   \xrightarrow{\epsilon\to0}
   -\frac{\beta}{g}.
\label{eq:(2.20)}
\end{equation}
在只减除极点项的 MS 方案中，这意味着
\begin{equation}
   Z_S=\left(1-\frac{\beta}{g}\frac{1}{\epsilon}\right)Z
   =1+O(g^4),
\label{eq:(2.21)}
\end{equation}
而式~\eqref{eq:(2.14)} 随后给出
\begin{equation}
   Z_M=-\frac{\beta}{4g}Z\frac{1}{\epsilon}
   =\frac{b_0}{4}g^2\frac{1}{\epsilon}+O(g^4).
\label{eq:(2.22)}
\end{equation}
于是我们看到，式~\eqref{eq:(2.10)} 与~\eqref{eq:(2.11)} 中的所有重正化常数
$Z_T$、$Z_M$ 与~$Z_S$ 最终都能用式~\eqref{eq:(2.7)} 中的规范耦合重正化常数
$Z$ 表示出来。
